\scanpage{133}
% archive foliation: 132
\deleted{\unsure{(iv)}}

$S$ préschém. loc. noeth.

$F$ foncteur
\diagram{\begin{tikzcd}(\text{Sch})_{/S}\arrow[r] & (\text{Ens})\end{tikzcd}}

(i) $F$ compatible aux limites inductives d'anneaux

(ii) $F$ compatible aux descentes \unsure{fidèlement} plats qua. cpts. \inserted{et localisation}

[(iii) Si $A$ anneau local noeth. \uline{complet} sur $S$, corps résiduel $k$, alors $F(A)\xrightarrow{\sim}F(k)$,

(iv) $F$ séparé.

(v) Pour \unsure{tt} corps $k$ sur $S$, $F(k)$ est \uline{fini}, et si $K'$ est une ext. alg. close de $k$, \deleted{\unsure{plat|$k$}} alg. clos, $F(k)\unsure{-|\to}F(k')$ est bijectif.

Alors $F$ représentable par un \deleted{foncteur} objet étale \deleted{et qu.} séparé qu. fini sur $S$.

\deleted{\unsure{Recommençons}}

Supposons $S$ le spectre d'un anneau local complet $A$, \inserted{à corps résiduel alg. clos.} \deleted{et que} de dim $\unsure{n|1}$, et qu'on sache que sur $S_1=S-\unsure{y|s}$, le foncteur $F$ soit représentable par $M_1$. \deleted{Considérons les \unsure{points}} Soient $\xi_i^0\deleted{\unsure{\in F_0/k}}$ \inserted{$(i\in I)$}, éléments de $F(k)$, ils proviennent d'éléments $\xi_i\in F(A)$ bien déterminés, considérons leurs restrictions $\xi_i'=\xi_i|S_1\in F(S_1)$, ils correspondent à des
